\documentclass{beamer}

\usepackage{hyperref}
\usepackage[utf8]{inputenc}
\usepackage[T1]{fontenc}
\usepackage[spanish]{babel}

% --- Paquetes para diseño ---
\usepackage{latexsym,amsmath,xcolor,multicol,booktabs,calligra}
\usepackage{graphicx,listings,stackengine,tikz}
\usetikzlibrary{arrows.meta,positioning,calc}

% --- Definición de Cajas Personalizadas ---
\newenvironment{cajaazul}[1]{%
    \setbeamercolor{block title}{bg=cyan!70!black, fg=white}
    \setbeamercolor{block body}{bg=cyan!10!white, fg=black}
    \begin{block}{#1}}{\end{block}}

\newenvironment{cajaverde}[1]{%
    \setbeamercolor{block title}{bg=teal!80!black, fg=white}
    \setbeamercolor{block body}{bg=teal!10!white, fg=black}
    \begin{block}{#1}}{\end{block}}

\newenvironment{cajaranja}[1]{%
    \setbeamercolor{block title}{bg=orange!90!black, fg=white}
    \setbeamercolor{block body}{bg=orange!10!white, fg=black}
    \begin{block}{#1}}{\end{block}}

% --- DATOS DEL PROYECTO ---
\author{Juan Pedro García Sanz \and Pablo Revuelto de Miguel \and \\ Adolfo Peña Marín}
\title{Robochess}
\subtitle{Dos brazos robóticos jugando al ajedrez de forma autónoma}
\institute{Robótica Inteligente \\ Universidad Loyola Andalucía}

% El logo se añade aquí, justo encima de la fecha con un salto de línea
\date{\includegraphics[height=1cm]{logo.png} \\[0.2cm] \today}

% --- TU TEMA PERSONALIZADO ---
\usepackage{CNU}
\usefonttheme{professionalfonts}
\usepackage{lmodern}

% --- Colores para código ---
\definecolor{deepblue}{rgb}{0,0,0.5}
\definecolor{deepred}{rgb}{0.6,0,0}
\definecolor{deepgreen}{rgb}{0,0.5,0}
\definecolor{halfgray}{gray}{0.55}

\newcommand{\completar}[1]{\textcolor{halfgray}{\textit{[#1]}}}

\begin{document}

% --- PORTADA ---
\begin{frame}
    \titlepage
\end{frame}

% ============================================================
% ÍNDICE
% ============================================================
\begin{frame}{Índice}
    \tableofcontents[hideallsubsections]
\end{frame}

% ============================================================
% 1. INTRODUCCIÓN Y ARQUITECTURA
% ============================================================
\section{Introducción y Arquitectura}

\subsection{Objetivo}

\begin{frame}{Qué es Robochess}

\begin{cajaazul}{La idea}
Dos manipuladores \textbf{ABB IRB120} enfrentados juegan una partida de
ajedrez de forma autónoma en simulación: le dictas la jugada por voz, el
sistema la valida y el brazo la ejecuta sobre el tablero.
\end{cajaazul}

\begin{cajaverde}{Por qué tiene su gracia}
No es mover un brazo aislado. Hay que encajar en un mismo flujo cuatro mundos
que normalmente van por separado: \textbf{percepción} (visión), \textbf{interfaz
por voz}, \textbf{razonamiento} sobre las reglas y \textbf{planificación de
movimiento}. El reto está en que todo eso hable entre sí.
\end{cajaverde}

\end{frame}

\subsection{Arquitectura}

\begin{frame}{Arquitectura: nodos que se hablan}

\begin{columns}[c]
\begin{column}{0.52\textwidth}
\centering
\begin{tikzpicture}[
    font=\small,
    nodo/.style={
        rounded corners, draw, thick, align=center,
        minimum width=2.6cm, minimum height=0.95cm, fill=cyan!10!white,
        draw=cyan!70!black},
    sensor/.style={nodo, fill=teal!12!white, draw=teal!80!black},
    actua/.style={nodo, fill=orange!12!white, draw=orange!90!black},
    flecha/.style={-{Stealth[length=2.5mm]}, thick, draw=gray!70!black},
]
\node[sensor] (voz)   {\textbf{Voz}\\ \scriptsize micrófono};
\node[sensor] (cam)   [right=0.9cm of voz] {\textbf{Cámara}\\ \scriptsize tablero};
\node[nodo]   (brain) [below=0.8cm of $(voz)!0.5!(cam)$]
    {\textbf{Cerebro}\\ \scriptsize decide y valida};
\node[nodo]   (mov)   [below=0.7cm of brain]
    {\textbf{Movimiento}\\ \scriptsize planifica};
\node[actua]  (sim)   [below=0.7cm of mov]
    {\textbf{Simulación}\\ \scriptsize ejecuta};

\draw[flecha] (voz) -- (brain.north);
\draw[flecha] (cam) -- (brain.north);
\draw[flecha] (brain) -- (mov);
\draw[flecha] (mov) -- (sim);
\draw[flecha, dashed, draw=teal!70!black]
    (sim.east) to[out=0, in=0]
    node[right=1pt, midway, font=\scriptsize, text=teal!50!black, align=left]
    {valida} (cam.east);
\end{tikzpicture}
\end{column}

\begin{column}{0.48\textwidth}
\begin{cajaverde}{Todo sobre ROS\,2 Jazzy}
Cada bloque es uno o varios \textbf{nodos} independientes que se comunican por:
\begin{itemize}\setlength{\itemsep}{2pt}
    \item \textbf{Topics}: flujo continuo (la imagen de la cámara).
    \item \textbf{Servicios}: petición–respuesta (interpretar una frase).
    \item \textbf{Acciones}: tareas largas con \textit{feedback} (ejecutar la jugada).
\end{itemize}
Así cada parte se desarrolla y se prueba por separado.
\end{cajaverde}
\end{column}
\end{columns}

\end{frame}

% ============================================================
% 2. PERCEPCIÓN: VISIÓN Y VOZ
% ============================================================
\section{Percepción: Visión y Voz}

\subsection{Visión}

\begin{frame}{Visión: del píxel a la casilla}

\begin{cajaazul}{Detección de piezas}
Una cámara cenital observa el tablero. Un detector \textbf{YOLOv8} localiza
cada pieza y la clasifica en \textbf{12 clases} (6 tipos $\times$ 2 colores),
devolviendo una caja por pieza.
\end{cajaazul}

\begin{cajaverde}{De la caja a la coordenada}
El centro de cada caja pasa por una \textbf{homografía} cenital ---calibrada
con la altura y el campo de visión de la cámara--- que convierte píxeles en
coordenadas del tablero, y de ahí a casilla (a1--h8). El resultado es un
estado del tablero ``observado''.

\smallskip
\textit{Rol:} la visión \textbf{verifica}, no manda. Contrasta lo que ve con
el estado interno; la autoridad es siempre el motor de ajedrez.
\end{cajaverde}

\end{frame}

\subsection{Voz}

\begin{frame}{Voz: del audio a una jugada legal}

\begin{cajaazul}{El recorrido}
Detección de actividad de voz (\textbf{VAD}) que corta cuando dejas de hablar
$\rightarrow$ \textbf{Whisper} transcribe a texto $\rightarrow$ un \textbf{LLM}
traduce la frase a notación de ajedrez (\texttt{e2e4}).
\end{cajaazul}
\begin{cajaverde}{Una decisión que nos salvó}
El parser solo puede devolver un movimiento que esté en la lista de jugadas
\textbf{legales}: nunca se inventa nada. Y si el LLM no responde o la conexión
falla, un parser por \textbf{expresiones regulares} hace de red de seguridad.
Robusto incluso con transcripciones ruidosas.
\end{cajaverde}

\end{frame}

% ============================================================
% 3. DECISIÓN Y MOVIMIENTO
% ============================================================
\section{Decisión y Movimiento}

\subsection{Cerebro}

\begin{frame}{El cerebro: las reglas mandan}

\begin{cajaazul}{Estado y validación}
El motor (\textbf{python-chess}) mantiene el estado autoritativo de la partida
en formato \textbf{FEN}. Para cada jugada propuesta comprueba legalidad, turno
correcto y detecta \textbf{jaque, mate y tablas}.
\end{cajaazul}

\begin{cajaverde}{El director de orquesta}
Si la jugada es válida, la empaqueta como objetivo y se la pasa al subsistema
de movimiento a través de una \textbf{acción} ROS\,2. Cuando el brazo termina,
actualiza el estado y cede el turno al otro bando.
\end{cajaverde}

\end{frame}

\subsection{Movimiento}

\begin{frame}{Movimiento: planificar sin chocar}

\begin{cajaazul}{MoveIt\,2}
Cada brazo tiene 6 grados de libertad. \textbf{MoveIt\,2} resuelve la
\textbf{cinemática inversa} y busca con \textbf{OMPL} una trayectoria libre de
colisiones hasta la pieza y hasta su destino, evitando tablero, piezas y el
otro brazo.
\end{cajaazul}

\begin{cajaverde}{Pick \& place por fases}
Aproximar, agarrar, elevar, transportar, soltar y retirarse. El agarre es
\textbf{magnético} para simplificar la manipulación de piezas pequeñas.
Detalle de diseño: reutilizamos un único planificador en marcha en lugar de
levantar un segundo, evitando duplicar trabajo.
\end{cajaverde}

\end{frame}

% ============================================================
% 4. INTEGRACIÓN Y RESULTADOS
% ============================================================
\section{Integración y Resultados}

\subsection{Simulación}

\begin{frame}{El escenario simulado}

\begin{cajaazul}{Gazebo Harmonic}
La escena: mesa, tablero, los dos IRB120 (blancas y negras) y la cámara
cenital, con física realista de peso y contactos.
\end{cajaazul}

\begin{cajaverde}{Control de los brazos}
Las articulaciones se mueven con \textbf{ros2\_control} a través del puente
ROS--Gazebo, con un controlador de trayectoria por brazo. Es la misma
infraestructura que se usaría sobre un robot real.
\end{cajaverde}

\end{frame}

\subsection{Resultados}

\begin{frame}{Qué funciona y qué nos costó}

\begin{cajaazul}{Lo que funciona hoy}
\begin{itemize}\setlength{\itemsep}{2pt}
    \item Partida dictada por voz, de principio a fin.
    \item Validación de reglas en tiempo real.
    \item \textit{Pick \& place} ejecutado por los dos brazos.
    \item Visión operativa como validador del tablero.
\end{itemize}
\end{cajaazul}

\begin{cajaranja}{Los retos que de verdad costaron}
\begin{itemize}\setlength{\itemsep}{2pt}
    \item Capturar el micrófono bajo WSL2 (audio que entraba en silencio).
    \item Que la pieza siguiera al gripper durante el transporte.
    \item Sobrevivir a transcripciones ruidosas (de ahí el \textit{fallback}).
\end{itemize}
\end{cajaranja}

\end{frame}

\subsection{Conclusiones}

\begin{frame}{Conclusiones}

\begin{cajaazul}{Lo que nos llevamos}
Integrar IA ---visión, voz y lenguaje--- con robótica ---planificación y
control--- sobre ROS\,2, y que todo funcione junto. La arquitectura modular
fue clave: cada pieza es sustituible sin tocar el resto.
\end{cajaazul}

\begin{cajaverde}{Siguientes pasos}
\begin{itemize}\setlength{\itemsep}{2pt}
    \item Un motor que juegue solo (p.\,ej. Stockfish) para robot contra humano.
    \item Saltar de la simulación al hardware real (calibración con marcadores).
    \item Afinar el detector de visión y los tiempos de planificación.
\end{itemize}
\end{cajaverde}

\end{frame}

\begin{frame}

\begin{center}
\Large \textbf{¿Preguntas?}
\end{center}

\end{frame}

\end{document}
